% Source: https://artofproblemsolving.com/community/c3h2079716_obscure_nt_trick
% Stable identity: aops-topic-2079716-q01
(a) By unique factorisation, a positive integer \(d\) divides \(n\) exactly when \(d=p_1^{a_1}p_2^{a_2}\cdots p_k^{a_k}\) with \(0\le a_i\le e_i\) for each \(i\). (Any prime factor of \(d\) is a prime factor of \(n\), and its exponent in \(d\) cannot exceed its exponent in \(n\); conversely every such product divides \(n\).) Different exponent choices give different divisors, again by unique factorisation. There are \(e_i+1\) choices for each \(a_i\), so by the multiplication principle the number of divisors is \((e_1+1)\cdots(e_k+1)\).

(b) \(10!=2^8\cdot 3^4\cdot 5^2\cdot 7\) (the exponent of \(2\) is \(5+2+1=8\), of \(3\) is \(3+1=4\), of \(5\) is \(2\)). So \(10!\) has
\[
(8+1)(4+1)(2+1)(1+1)=\boxed{270}
\]
positive divisors.
